% =============================================================================
%  CHAPITRE 1 : CADRE CONCEPTUEL ET ÉTAT DE L'ART
% =============================================================================
\entete{Chapitre 1 : Cadre conceptuel et état de l'art}
\chapter{Cadre conceptuel et état de l'art}
\label{chap:cadre}

\introchap

Ce chapitre pose les bases du travail. Il commence par définir les notions
employées tout au long du rapport. Cette clarification n'est pas une formalité :
des mots comme « immersion », « visite virtuelle » ou « réalité augmentée » sont
utilisés avec des sens très différents selon les auteurs et selon les produits
commerciaux, au point qu'ils finissent par ne plus rien désigner de précis. Le
chapitre examine ensuite ce qui existe déjà, d'abord dans les musées
camerounais, puis à l'échelle internationale, afin de situer notre contribution.

% =============================================================================
\section{Cadre conceptuel}
\label{sec:conceptuel}

\subsection{Patrimoine culturel et valorisation numérique}

Le patrimoine culturel désigne l'ensemble des biens, matériels et immatériels,
qu'une communauté reconnaît comme porteurs de son identité et qu'elle entend
transmettre ; la Convention de l'UNESCO (1972) distingue le patrimoine
\emph{matériel}, monuments, ensembles architecturaux et objets, du patrimoine
\emph{immatériel}, savoir-faire, rites et traditions orales. À la Fondation Jean
Félicien Gacha, les deux dimensions coexistent : les objets exposés ne prennent
leur sens que rapportés aux lignées, aux fonctions et aux usages dont ils
procèdent.

La \emph{valorisation numérique} consiste à produire, organiser et diffuser des
représentations numériques de ces biens afin d'en élargir l'accès et d'en
assurer la conservation documentaire. Elle se distingue de la simple
\emph{numérisation}, opération technique convertissant un objet en fichier :
numériser, c'est photographier ; valoriser, c'est rendre intelligible, relier et
transmettre. La distinction est essentielle ici, car le problème posé par la
Fondation n'était pas un manque de photographies, mais un manque de mise en
relation et de récit.

\subsection{De la photographie à la visite immersive}

Une visite virtuelle est un parcours numérique permettant d'explorer un lieu
sans s'y rendre, et trois formes en coexistent. \textbf{La galerie de
photographies}, suite d'images fixes éventuellement légendées, est la plus
répandue, la moins coûteuse et la moins immersive : c'est celle que la Fondation
pratiquait. \textbf{Le parcours panoramique} enchaîne des prises de vue à
360 degrés reliées par des points de passage, l'utilisateur pivotant librement
et percevant les volumes de la salle. \textbf{La scène tridimensionnelle},
enfin, reconstruit le lieu et laisse l'utilisateur y circuler comme il
marcherait : la plus proche d'une visite réelle, mais exigeant une
reconstruction lourde et une machine puissante côté visiteur.

Un mot doit être précisé, car il est souvent employé à tort. Il convient de
distinguer l'\emph{immersion}, propriété technique du dispositif (champ de
vision, temps de réaction, qualité d'image), de la \emph{présence}, sentiment
d'« être là » éprouvé par l'utilisateur (Slater et Wilbur, 1997) : un dispositif
techniquement immersif ne produit pas automatiquement ce sentiment, auquel la
continuité du déplacement et la qualité du récit contribuent autant que la
technique. Ce constat oriente nos choix : plutôt que l'immersion maximale au
prix d'un équipement inaccessible, nous avons cherché le meilleur sentiment de
présence possible sur le matériel dont dispose réellement le public visé, un
téléphone ou un ordinateur ordinaire.

\subsection{La réalité augmentée}

La réalité augmentée superpose des éléments virtuels, images ou objets en trois
dimensions, sur le monde réel et en temps réel, au travers de l'écran d'un
téléphone ou d'une tablette : l'appareil analyse ce que voit sa caméra, y détecte
les surfaces, puis y ancre l'élément virtuel de telle sorte qu'il paraisse posé
dans la pièce et y reste quand l'utilisateur se déplace autour de lui. Milgram
et Kishino (1994) la situent sur un continuum allant du monde réel au monde
entièrement virtuel : la réalité virtuelle remplace l'environnement de
l'utilisateur, la réalité augmentée s'y ajoute.

Trois choses souvent réunies sous le mot « 3D » doivent être distinguées, car
elles répondent à des besoins différents et ne portent pas sur les mêmes sujets.
\textbf{Faire tourner une œuvre à l'écran} la montre sous tous ses angles, mais
elle reste enfermée dans l'écran. \textbf{Parcourir un lieu} en panoramas
enchaînés donne la sensation de s'y déplacer. \textbf{Poser un espace chez soi
en réalité augmentée}, enfin, le fait sortir de l'écran pour apparaître à sa
taille véritable : c'est le seul dispositif qui réponde à la question « quelle
taille cela fait-il ? », et c'est pourquoi il s'applique aux habitations et aux
espaces plutôt qu'aux objets. Le dispositif construit emploie les trois, et se
situe volontairement du côté de la réalité augmentée \emph{sans équipement}, au
moyen de la seule caméra du téléphone : exiger un casque reviendrait à recréer,
sous une autre forme, la barrière d'accès que le projet entend supprimer.

\subsection{Acquisition et modélisation en trois dimensions}

Trois voies permettent d'obtenir le modèle tridimensionnel d'un objet ou d'un
lieu patrimonial. \textbf{La modélisation manuelle}, réalisée par un
infographiste, offre un contrôle total sur le résultat mais coûte cher et ne
garantit pas la fidélité aux dimensions réelles. \textbf{La photogrammétrie}
reconstruit la géométrie à partir de photographies prises sous des angles
différents, un même point vu depuis plusieurs positions permettant de calculer
sa position dans l'espace ; elle demande peu de matériel, mais beaucoup de
méthode dans la prise de vue et de la puissance de calcul. \textbf{La télémétrie
par laser}, ou LiDAR, mesure directement les distances : depuis que ce capteur
équipe les téléphones haut de gamme, elle fournit en quelques minutes une
géométrie correctement dimensionnée.

Les modèles obtenus doivent ensuite être diffusés. Le format retenu par
l'industrie du web est le glTF, avec sa variante compacte GLB, que le consortium
Khronos (2021) présente comme le format d'échange universel de la 3D ; les
appareils d'Apple, eux, n'acceptent que l'USDZ. Cette dualité n'est pas un
détail technique : produire un seul des deux formats revient à priver de la
fonction la moitié des visiteurs.

\subsection{Progiciel de gestion intégré et cloisonnement des données}

Un progiciel de gestion intégré, ou ERP, centralise les processus de gestion
d'une organisation, inventaire des collections, espaces, billetterie, boutique
et événements, autour d'une base de données unique. Appliqué à une institution
patrimoniale, il remplace les cahiers et les fichiers dispersés par un système
consultable et sauvegardé.

Le modèle du logiciel en tant que service consiste à faire fonctionner une seule
instance de l'application pour plusieurs organisations clientes. On parle alors
d'architecture multi-organisations. Chong et Carraro (2006) en distinguent trois
mises en œuvre : une base de données par organisation, un schéma par
organisation, ou une base commune avec cloisonnement ligne par ligne. La
troisième est la plus économique et la plus délicate. Dans une architecture en
nuage de ce type, l'isolation des données repose sur la \emph{sécurité au niveau
de la ligne}, mécanisme du système de gestion de base de données qui filtre les
lignes visibles selon l'identité de celui qui interroge. Ce mécanisme garantit
qu'une organisation ne peut consulter que ses propres données, et cette garantie
est appliquée par le moteur de la base lui-même, non par le programme qui
l'interroge. La distinction se révélera décisive au moment du diagnostic.

\subsection{Intelligence artificielle générative et génération augmentée par la recherche}

Un grand modèle de langue est un réseau de neurones entraîné à prédire la suite
d'un texte ; l'architecture qui domine aujourd'hui, le \emph{Transformer}, a été
introduite par Vaswani et al. (2017). Ces modèles produisent un texte fluide et
grammaticalement correct, mais ils ne disposent d'aucun mécanisme interne de
vérification des faits : ils peuvent énoncer avec assurance des informations
fausses, ce que la littérature nomme \emph{hallucination} (Ji et al., 2023).
Dans un contexte de médiation culturelle, où l'exactitude prime sur l'aisance,
ce défaut est rédhibitoire s'il n'est pas corrigé.

La correction passe par deux notions. La première est le \emph{plongement}, ou
\emph{embedding} : la représentation d'un texte sous forme de suite de nombres,
construite de telle manière que deux textes de sens voisin soient représentés
par deux suites proches. Cette représentation permet de chercher par le sens et
non par les mots : une question portant sur un « masque de danse » fait alors
remonter une notice décrivant un « masque cérémoniel », alors qu'aucun mot n'est
commun aux deux formulations. L'interrogation rapide de ces représentations
suppose un index adapté, la structure HNSW proposée par Malkov et Yashunin
(2018) étant aujourd'hui la plus employée.

La seconde notion est la \textbf{génération augmentée par la recherche}, en
anglais \emph{retrieval-augmented generation}, formalisée par Lewis et al.
(2020). Le principe consiste à ne pas demander au modèle de répondre de mémoire,
mais à lui fournir, au moment de la question, les documents pertinents extraits
d'une base maîtrisée, en lui imposant de fonder sa réponse sur eux seuls. Le
mécanisme se décompose en cinq temps :

\begin{enumerate}
  \item \textbf{l'ingestion} : les documents de référence sont collectés et
        découpés en fragments de taille homogène ;
  \item \textbf{l'indexation} : chaque fragment est converti en plongement, puis
        enregistré dans un index ;
  \item \textbf{la recherche} : la question est convertie de la même manière et
        comparée à l'index, qui renvoie les fragments les plus proches ;
  \item \textbf{l'augmentation} : ces fragments sont insérés dans l'instruction
        adressée au modèle, avec la consigne de n'affirmer que ce qui y figure ;
  \item \textbf{la génération} : le modèle rédige la réponse et cite ses
        sources ; s'il ne trouve pas l'information, il doit le dire.
\end{enumerate}

L'intérêt de ce dispositif pour une institution patrimoniale est direct : il
procure la fluidité d'un guide et la fiabilité d'un catalogue, et il rend la
mise à jour immédiate, puisque modifier une notice suffit à modifier les
réponses, sans réentraîner quoi que ce soit.

\subsection{DevOps : industrialiser la production du logiciel}

Le DevOps est une manière d'organiser le travail qui vise à raccourcir le chemin
séparant une modification écrite par un développeur de sa mise à disposition des
utilisateurs. Son nom, contraction de \emph{development} et \emph{operations},
dit le projet : supprimer la frontière entre ceux qui écrivent le logiciel et
ceux qui le font fonctionner. Kim et al. (2016) en résument le principe par
trois exigences, fluidifier le flux vers la production, raccourcir les boucles
de retour et instaurer une culture d'apprentissage continu ; Bass, Weber et Zhu
(2015) en précisent les conséquences sur l'architecture des systèmes.

Deux pratiques, formalisées par Humble et Farley (2010), en constituent le cœur.
L'\textbf{intégration continue} veut que chaque modification déposée déclenche
automatiquement la compilation et une série de vérifications, de sorte qu'une
anomalie soit signalée en quelques minutes et non le jour de la mise en ligne.
Le \textbf{déploiement continu} veut qu'une modification acceptée soit publiée
selon une procédure rigoureusement identique à chaque fois, ce qui supprime les
erreurs de manipulation. Deux techniques les rendent possibles : le
\textbf{conteneur}, qui enferme l'application avec tout ce dont elle a besoin
pour fonctionner, de sorte que ce n'est plus le code seul qui voyage mais son
environnement complet ; et l'\textbf{infrastructure décrite par le code}, qui
consigne dans des fichiers versionnés l'ensemble des ressources à créer chez
l'hébergeur, qu'un outil tel que Terraform (HashiCorp, 2024) crée ensuite et
maintient conformes.

Ce recours au DevOps n'est pas ici un choix de confort. Deux mois de stage ne
laissent aucune marge pour reconstruire ce qui aurait été cassé sans qu'on s'en
aperçoive, et la Fondation devra un jour exploiter le système sans nous : une
mise en production qui repose sur une suite de gestes mémorisés par une seule
personne n'est pas transmissible, une mise en production décrite dans un fichier
l'est. La nature des pannes rencontrées le confirme : un modèle tridimensionnel
servi avec une étiquette de type erronée s'affiche parfaitement dans le
navigateur mais empêche toute session de réalité augmentée, anomalie qui ne se
voit pas à l'œil et se vérifie automatiquement, ou passe inaperçue jusqu'au jour
de la démonstration. Une équipe technique réduite, enfin, ne peut pas consacrer
une journée à chaque mise en ligne : automatiser la publication, c'est rendre
les mises à jour assez peu coûteuses pour qu'elles aient effectivement lieu.

\subsection{FinOps : maîtriser le coût de l'infrastructure en nuage}

L'hébergement en nuage transforme une dépense d'investissement en dépense de
fonctionnement : on n'achète plus un serveur, on paie ce que l'on consomme.
Cette souplesse a une contrepartie, car la dépense devient variable, diffuse et
facilement invisible, personne ne recevant de facture avant d'avoir consommé. La
discipline dite \emph{FinOps} (FinOps Foundation, 2023) vise à rendre ce coût
visible, à l'attribuer à ceux qui le provoquent et à l'optimiser en continu ;
elle procède en trois temps : \textbf{informer}, mesurer et répartir la dépense ;
\textbf{optimiser}, supprimer le gaspillage et dimensionner au juste besoin ;
\textbf{opérer}, inscrire ces contrôles dans le fonctionnement courant plutôt
que de les réserver aux audits.

Pour une institution à budget contraint, la question n'est pas de savoir si le
dispositif fonctionne, mais combien il coûtera une fois le stagiaire parti ; elle
a d'ailleurs été posée telle quelle lors des entretiens. Trois apports y
répondent. D'abord un chiffre plutôt qu'une intuition : les montants
d'infrastructure présentés dans ce rapport ne sont pas déduits d'un barème, ils
sont calculés service par service par \emph{Infracost} à partir des fichiers qui
décrivent l'infrastructure, et donc connus avant même que la dépense ne soit
engagée. Ensuite une architecture pensée pour la dépense : servir l'application
sous forme de fichiers statiques diffusés par un réseau de distribution, plutôt
que par un serveur allumé en permanence, fait tomber le coût récurrent à
quelques centaines de francs par mois, le seul poste réellement coûteux, le
calcul graphique de la photogrammétrie, étant ponctuel et s'interrompant de
lui-même hors usage. Enfin des garde-fous durables : étiqueter les ressources,
poser des alarmes budgétaires et arrêter automatiquement les environnements
d'essai ne coûte rien et évite la dérive silencieuse qui caractérise les
dépenses de nuage.

% =============================================================================
\section{État de l'art}
\label{sec:etatart}

\subsection{La situation des musées camerounais}

L'état de l'art doit d'abord être dressé là où se situe le problème, c'est-à-dire
au Cameroun. Le pays compte une trentaine d'établissements muséaux, répartis
entre musées nationaux, musées de chefferie et musées privés ou confessionnels.
Le tableau~\ref{tab:museescmr} présente les principaux d'entre eux et le mode de
présentation qu'ils offrent aujourd'hui à un public distant.

\begin{table}[!ht]
\centering
\caption{Principaux musées camerounais et mode de présentation à distance}
\label{tab:museescmr}
\small
\begin{tabularx}{\textwidth}{|L{4.6cm}|L{2.6cm}|X|}
\hline
\rowcolor{keycetete}
\thead{Établissement} & \thead{Localisation} & \thead{Présentation offerte au public distant} \\ \hline
Musée national du Cameroun & Yaoundé & Photographies et notices ; essai de visualisation 3D sur une exposition temporaire \\ \hline
Musée des civilisations & Dschang & Photographies et notices \\ \hline
Musée du palais royal & Foumban & Photographies et notices \\ \hline
Musée de la chefferie & Bandjoun & Photographies et notices \\ \hline
Musée de la chefferie & Baham & Photographies et notices \\ \hline
Musée de la chefferie & Babungo & Photographies et notices \\ \hline
Musée Blackitude & Yaoundé & Photographies et notices \\ \hline
Musée maritime & Douala & Photographies et notices \\ \hline
Fondation Jean Félicien Gacha & Bangoulap & Photographies et publications sur les réseaux sociaux \\ \hline
\end{tabularx}
\source{Relevé effectué par nos soins en juillet 2026 à partir des pages
publiques de ces établissements et de Keumoe et Blot (2023).}
\end{table}


Ce relevé appelle trois constats.

\begin{enumerate}
  \item \textbf{Aucun musée camerounais ne propose de visite immersive à ses
        visiteurs.} Là où une présentation en ligne existe, elle se limite à
        des photographies accompagnées de notices descriptives. Cette pratique,
        que l'on peut qualifier de photogrammétrie statique, consiste à
        photographier l'objet et à le décrire : elle documente, mais elle ne
        restitue ni le volume, ni l'échelle, ni le lieu.

  \item \textbf{Le seul travail de visualisation en trois dimensions documenté
        est expérimental.} Keumoe et Blot (2023) rapportent une initiative
        conduite au Musée national du Cameroun autour de l'exposition
        temporaire \emph{Contact(s) Zone}. Il s'agissait d'un projet pilote mené
        par des chercheurs et non d'une commande de l'établissement ; la
        modélisation a été réalisée sous Blender et restituée avec Unreal
        Engine~5, la photogrammétrie, pourtant plus fidèle, ayant été écartée
        faute de moyens. Le dispositif portait sur une exposition temporaire et
        n'a pas été mis en service pour le public.

  \item \textbf{Les obstacles sont connus et convergents.} Les mêmes auteurs
        relèvent l'insuffisance des moyens techniques, l'absence de formation du
        personnel aux sciences du patrimoine, le manque de plans architecturaux
        fiables, la précarité financière et la faible synergie entre
        institutions culturelles du pays.
\end{enumerate}

Il existe donc, au Cameroun, un écart entre la richesse du patrimoine conservé
et les moyens de le faire connaître au-delà des murs. C'est cet écart que notre
travail entend réduire, en montrant qu'un dispositif immersif complet peut être
construit et exploité avec les moyens d'une institution ordinaire.

\subsection{Les dispositifs proposés à l'international}

Hors du Cameroun, le marché de la visite virtuelle s'est structuré autour de
quelques familles de solutions, dont les caractéristiques sont résumées au
tableau~\ref{tab:comparvv}.

\begin{table}[!ht]
\centering
\caption{Comparaison des principales familles de solutions de visite virtuelle}
\label{tab:comparvv}
\small
\begin{tabularx}{\textwidth}{|L{3.1cm}|X|L{3.0cm}|L{2.6cm}|}
\hline
\rowcolor{keycetete}
\thead{Solution} & \thead{Principe et apports} & \thead{Limites relevées} & \thead{Modèle économique} \\ \hline
Google Arts \& Culture &
Numérisation à très haute définition et parcours en vue immersive ; audience
mondiale considérable &
Sélection à l'initiative de l'opérateur ; l'institution ne maîtrise ni sa
présentation ni ses données ; aucun outil de gestion &
Gratuit, sur invitation \\ \hline
Matterport &
Capture tridimensionnelle d'un lieu par caméra dédiée ; rendu de grande qualité &
Matériel spécifique coûteux ; hébergement des données chez l'éditeur ; ni
assistance conversationnelle ni gestion documentaire &
Abonnement mensuel par espace \\ \hline
Kuula, Pano2VR &
Assemblage de panoramas à 360 degrés avec points d'intérêt ; prise en main
rapide et coût faible &
Pas de base documentaire, pas de réalité augmentée, pas de gestion de
l'institution ; parcours isolé du reste du site &
Abonnement, palier gratuit \\ \hline
Développement sur mesure &
Adapté au besoin exact de l'institution &
Coût de conception élevé, non réutilisable pour une autre structure, maintenance
à la charge du commanditaire &
Prestation ponctuelle \\ \hline
\end{tabularx}
\source{Élaboré par nos soins à partir de la documentation publique des éditeurs,
consultée en juin 2026.}
\end{table}


Une constante ressort de cette comparaison : ces solutions traitent la visite
virtuelle comme un produit isolé, détaché de la gestion de l'institution et de
sa documentation. Aucune ne relie le parcours immersif à un inventaire, à une
billetterie, à une base de connaissances interrogeable ou à d'autres
institutions. C'est précisément cette articulation qui fait défaut.

\subsection{Les collections ouvertes et leurs interfaces de programmation}

Plusieurs grands musées ont ouvert leurs catalogues au moyen d'interfaces de
programmation publiques. Ces sources constituent la matière première du
rapprochement documentaire recherché. Neuf d'entre elles ont été examinées.
Cinq sont interrogeables sans clé d'accès et ont été retenues : le Metropolitan
Museum of Art, dont la couverture africaine est la meilleure et qui a servi de
source principale ; l'Art Institute of Chicago ; le Cleveland Museum of Art,
dont les notices fournissent les numéros d'inventaire ; le Victoria and Albert
Museum, riche en arts appliqués ; et Wikidata, qui relie objets, cultures et
lieux sous forme de données structurées. Quatre autres, Europeana, le
Rijksmuseum, la Smithsonian Institution et les Harvard Art Museums, exigent une
clé que nous n'avons pu obtenir et ont dû être écartées.

Le choix des sources n'a donc pas été dicté par leur richesse, mais par leur
accessibilité effective dans les conditions du stage. Cette exclusion constitue
une limite explicite de l'étude et non un oubli.

\subsection{Les assistants conversationnels en contexte muséal}

Les premiers agents conversationnels muséaux reposaient sur des arbres de
questions prédéfinies : robustes mais rigides, ils ne répondaient qu'aux
questions prévues. L'arrivée des grands modèles de langue a permis une
conversation naturelle, au prix du risque d'invention évoqué plus haut. La
littérature récente converge vers l'ancrage documentaire comme réponse à ce
risque et formule trois exigences : la \textbf{traçabilité}, chaque affirmation
devant pouvoir être rapportée à un document identifiable ; l'\textbf{aveu
d'ignorance}, l'assistant devant pouvoir répondre qu'il ne sait pas, ce qui
suppose que la consigne l'y autorise ; et le \textbf{périmètre}, l'assistant
devant refuser poliment les questions hors sujet plutôt que d'y répondre
approximativement. Ces trois exigences ont directement structuré la conception
de notre assistant.
% =============================================================================
\section{Cadre juridique et éthique}
\label{sec:juridique}

Diffuser un patrimoine en ligne n'est pas un simple geste technique : cela pose
la question de savoir à qui appartiennent les images, et qui décide de ce qui
peut être montré.

\textbf{À qui appartiennent les images.} Il faut distinguer l'objet de sa
photographie. La plupart des pièces d'une chefferie sont anciennes et sans
auteur connu : personne ne détient de droit sur elles. Mais la photographie de
cet objet, et plus encore son modèle en trois dimensions, sont des créations
nouvelles, et celui qui les produit détient un droit dessus. Autrement dit,
photographier une œuvre libre de droits produit un fichier qui, lui, ne l'est
pas.

La règle retenue pour la plateforme est simple : \textbf{les images et les
modèles publiés sur le site d'une institution appartiennent à cette
institution}. Le visiteur peut les consulter librement, les faire tourner, les
regarder autant qu'il le souhaite, mais il n'est pas autorisé à les
réutiliser — ni à les télécharger pour son propre usage, ni à les republier
ailleurs, ni à s'en servir à des fins commerciales. Cette règle figure dans les
conditions du site, et le système conserve pour chaque média la trace de
l'institution qui l'a versé, afin qu'on puisse toujours dire à qui il
appartient.

\textbf{Ce que demandent les collections extérieures.} Le rapprochement
documentaire décrit au chapitre~3 interroge des musées étrangers qui ne
diffusent pas tous aux mêmes conditions : les uns laissent tout usage libre, les
autres l'interdisent à des fins commerciales, d'autres encore exigent d'être
cités. Pour respecter ces règles sans avoir à les traiter une par une, la
plateforme ne conserve jamais l'image d'un musée étranger : elle enregistre
seulement la notice et son adresse, et affiche à côté de chaque correspondance
le nom du musée détenteur et le numéro d'inventaire. Le lecteur peut ainsi
remonter à la source, ce qui est aussi une garantie de sérieux.

\textbf{Qui décide de ce qui est montré.} Une question dépasse le droit.
Certaines pièces ont une fonction rituelle, et la communauté qui les a produites
peut estimer qu'elles ne doivent pas être exposées à tous les regards. Aucune
loi ne tranche cela à sa place. C'est pourquoi rien n'est publié
automatiquement : un musée, une salle ou une œuvre reste invisible tant qu'une
personne de l'institution n'a pas décidé de la publier, et les rapprochements
avec d'autres musées ne s'affichent qu'après validation par le conservateur. Ces
deux mécanismes sont présentés plus loin comme des choix d'interface ; ce sont
d'abord des garanties éthiques, car ils laissent la décision à l'institution et
non au programme.

Sur la restitution des œuvres africaines, enfin, Sarr et Savoy (2018) ont posé
les termes du débat. Notre travail ne s'y substitue pas : montrer où se trouvent
les œuvres apparentées n'est pas les restituer, et le dispositif doit être
compris comme un outil de connaissance.

\textbf{Les données des visiteurs.} Créer un compte et recevoir des campagnes
d'information suppose de traiter des données personnelles. Trois règles s'y
appliquent : ces données ne servent qu'à la relation entre l'institution et son
adhérent ; le consentement est demandé explicitement et chaque message comporte
un lien de désinscription ; et la liste des destinataires est calculée par la
base de données, non par le programme qui rédige l'envoi. Ce dernier point
rejoint un enseignement plus général de ce travail, développé au chapitre~4 :
une garantie portée par la donnée résiste à l'erreur de programmation, une
garantie portée par le code n'y résiste pas.





\subsection{Ce que notre travail apporte}

Au terme de cet examen, trois manques se dégagent, qui définissent l'espace de
notre contribution.

\begin{enumerate}
  \item \textbf{L'absence de dispositif immersif dans les musées camerounais.}
        Aucun établissement ne propose aujourd'hui à ses visiteurs distants
        autre chose que des photographies légendées, et le seul travail de
        modélisation documenté est resté expérimental. Nous proposons un
        dispositif complet, en service, construit avec les moyens d'une
        institution ordinaire.

  \item \textbf{L'absence d'articulation.} Les solutions existantes séparent la
        visite immersive, la gestion de l'institution et l'assistance au
        visiteur. Nous les réunissons dans un ensemble où l'objet du parcours
        immersif est le même que celui de l'inventaire, de la boutique et de la
        base de connaissances de l'assistant.

  \item \textbf{L'absence de mise en relation entre institutions et de modèle
        reproductible.} Aucune solution examinée ne cherche, pour une œuvre
        donnée, ses œuvres apparentées ailleurs dans le monde, et aucune ne se
        transpose d'une institution à l'autre sans être reconstruite. Nous
        proposons un dispositif de rapprochement documentaire assorti d'une
        validation humaine, et une architecture ouverte dans laquelle l'accueil
        d'une nouvelle structure ne demande aucun développement.
\end{enumerate}

\conclchap

Ce chapitre a précisé le vocabulaire du rapport et situé le travail dans son
paysage. Les briques techniques nécessaires sont individuellement disponibles et
éprouvées, mais aucun musée camerounais n'en dispose aujourd'hui pour son public
distant. Ce qui manque n'est donc pas la technologie, mais son
\emph{assemblage} au service d'une institution à moyens limités. C'est cet
assemblage qui constitue l'objet de notre travail, et le chapitre suivant expose
la démarche adoptée pour le mener.
